\documentclass[11pt]{article}

\usepackage{amsmath,amssymb,amsthm}
\usepackage[margin=1.2in]{geometry}
\usepackage[colorlinks=true,linkcolor=blue,citecolor=blue]{hyperref}

\newtheorem{theorem}{Theorem}
\newtheorem{lemma}[theorem]{Lemma}
\newtheorem{proposition}[theorem]{Proposition}
\newtheorem{corollary}[theorem]{Corollary}
\theoremstyle{definition}
\newtheorem{definition}[theorem]{Definition}

\newcommand{\F}{\mathbb{F}}
\newcommand{\Z}{\mathbb{Z}}
\newcommand{\Q}{\mathbb{Q}}
\newcommand{\Clo}{\mathrm{Clo}}

\title{One-operation completeness for polynomial rings:\\
a characteristic-two dichotomy\\[1ex] {\large\normalfont --- condensed version ---}}
\author{John Thompson}
\date{June 2026}

\begin{document}
\maketitle

\begin{abstract}
Call $p \in R[x,y]$ \emph{complete} if every bivariate polynomial over
$R$ is an iterated composition of $p$, the variables, and constants.
We prove: a complete binary polynomial exists over a commutative ring
$R$ if and only if $2$ is a unit of $R$; when it is, $x^2 - y$ is
complete, and when it is not, for every $p$ at least one of $x+y$,
$xy$ is inexpressible.
\end{abstract}

\section{Definitions and statements}\label{sec:prelim}

\begin{definition}\label{def:clone}
Let $R$ be a commutative ring and $p \in R[x,y]$.  The
\emph{substitution clone} $\Clo(p) \subseteq R[x,y]$ is the smallest
set containing $x$, $y$, and every constant $c \in R$, and closed
under $(f,g) \mapsto p(f,g)$ (substitution of $f$ for $x$ and $g$ for
$y$ in $p$).  Call $p$ \emph{complete} over $R$ if
$\Clo(p) = R[x,y]$.
\end{definition}

Two points of precision.  First, $\Clo(p)$ is not a clone of
operations on a set: its elements are \emph{formal} polynomials ---
it is the binary component of the subclone of the abstract clone
(Lawvere theory) of commutative $R$-algebras generated by $p$, the
projections, and the constants; equivalently, the subalgebra of the
magma $\bigl(R[x,y],\, (f,g) \mapsto p(f,g)\bigr)$ generated by
$\{x,y\} \cup R$.  Second, the formal/functional distinction has
teeth: $n := 1 + xy \in \F_2[x,y]$ induces NAND on $\{0,1\}$ --- a
single-operation analogue of the Sheffer stroke \cite{Sheffer} --- so
the \emph{function} clone it generates contains every Boolean
function, yet $D n = 1 \neq 0$ for the operator $D$ of
Section~\ref{sec:negative}, so Theorem~\ref{thm:F2} gives
$x + y \notin \Clo(n)$: no composition of $n$ equals the polynomial
$x+y$ on the nose.

\begin{theorem}[dichotomy]\label{thm:dichotomy}
Let $R$ be a commutative ring.  A complete binary polynomial over $R$
exists if and only if $2$ is a unit of $R$.
\end{theorem}

\begin{theorem}\label{thm:converse}
If $\tfrac12 \in R$ then $q(x,y) := x^2 - y$ is complete:
$\Clo(x^2 - y) = R[x,y]$.
\end{theorem}

\begin{theorem}\label{thm:F2}
For every $q \in \F_2[x,y]$, $\;x + y \notin \Clo(q)$ or
$xy \notin \Clo(q)$.
\end{theorem}

There is no hypothesis on $q$: reduction mod $2$ collapses degree
($2x^3 + xy \mapsto xy$), so the deduction of
Corollary~\ref{cor:Z} below requires Theorem~\ref{thm:F2} for
arbitrary $q$, degenerate cases included.

\begin{theorem}\label{thm:char2}
For every field $k$ of characteristic $2$ and every $q \in k[x,y]$,
at least one of $x+y$, $xy$ lies outside $\Clo(q)$.
\end{theorem}

\begin{corollary}\label{cor:Z}
No binary polynomial over $\Z$ is complete.
\end{corollary}

\begin{lemma}\label{lem:criterion}
$\Clo(p) = R[x,y]$ if and only if $x + y \in \Clo(p)$ and
$xy \in \Clo(p)$.
\end{lemma}

\begin{proof}
Substituting members of the clone into a member lands in the clone
(induction on the formula).  So if $x+y, xy$ are members, the clone
is closed under $+$ and $\times$; containing the variables and
constants, it is everything.
\end{proof}

\begin{lemma}[transfer]\label{lem:transfer}
A ring homomorphism $R \to S$ induces
$\pi : R[x,y] \to S[x,y]$ with
$\pi\bigl(\Clo(p)\bigr) \subseteq \Clo(\pi p)$.
\end{lemma}

\begin{proof}
$\pi$ fixes variables, sends constants to constants, and commutes
with substitution; induct over the formula.
\end{proof}

By Lemmas~\ref{lem:criterion} and~\ref{lem:transfer} applied to
$\Z \to \F_2$, Theorem~\ref{thm:F2} implies Corollary~\ref{cor:Z}.

We use three facts about the formal partials
$\partial_x, \partial_y$ on $\F_2[x,y]$:
(1) $\partial(f^2) = 2f\,\partial f = 0$;
(2) $\partial_x^2 = \partial_y^2 = 0$, while
$D := \partial_x \partial_y$ need not vanish;
(3) squaring is a ring homomorphism, and a polynomial with all
exponents even is a square.
Univariate degrees follow the convention that constants ---
including $0$ --- have degree $0$, so $\deg(f^2) = 2\deg f$ without
exception.

\section{The positive half: $x^2 - y$}\label{sec:converse}

\begin{proof}[Proof of Theorem~\ref{thm:converse}]
By Lemma~\ref{lem:criterion} it suffices to reach $x+y$ and $xy$.
Everything flows from
\begin{equation}\label{eq:master-id}
  q\bigl(\alpha,\, q(\alpha', \gamma)\bigr)
  = \alpha^2 - \alpha'^2 + \gamma .
\end{equation}

\emph{Translations.}  With constants $\alpha = \frac{d+1}{2}$,
$\alpha' = \frac{d-1}{2}$, \eqref{eq:master-id} sends
$\gamma \mapsto \gamma + d$ for any $d \in R$ (the only step using
$\tfrac12$); in particular $x + d, y + d \in \Clo(q)$.

\emph{Slanted lines.}  $q(c, y) = c^2 - y$, then
$q(x, c^2 - y) = x^2 + y - c^2$, then
\[
  q\bigl(x + d,\; x^2 + y - c^2\bigr) = 2dx - y + (d^2 + c^2)
  \;\in\; \Clo(q),
\]
a line of slope $2d$ for every $d$, constant adjustable by
translations.

\emph{Linear engine.}  For $\alpha \in \Clo(q)$ and constant $k$,
\eqref{eq:master-id} with outer argument $\alpha + k$ and inner
argument $\alpha$ gives the increment
\[
  \gamma \;\longmapsto\; \gamma + 2k\,\alpha + k^2 .
\]
Increments along slanted lines of two distinct slopes span all
linear polynomials, so from $\gamma = 0$: every linear polynomial
lies in $\Clo(q)$, in particular $x + y$ and $x + \tfrac{y}{2}$.

\emph{Scaled squares.}  The engine with $\alpha = y^2 + \ell$ (in the
clone via $q(y, -\ell)$) adds $\lambda y^2 + \text{const}$ for any
$\lambda$; from $\gamma = x^2 = q(x, 0)$ we get
$x^2 + \tfrac14 y^2 \in \Clo(q)$.

\emph{Finish.}
\[
  q\Bigl(x + \tfrac{y}{2},\; x^2 + \tfrac14 y^2\Bigr)
  = \Bigl(x + \tfrac{y}{2}\Bigr)^2 - x^2 - \tfrac14 y^2
  = xy . \qedhere
\]
\end{proof}

\section{The negative half over $\F_2$}\label{sec:negative}

Fix $q \in \F_2[x,y]$.  Grouping the monomials of $q$ by the
parities of their exponents and using fact (3), $q$ has a unique
\emph{parity decomposition}
\begin{equation}\label{eq:parity}
  q \;=\; A_0^2 \;+\; x A_1^2 \;+\; y A_2^2 \;+\; xy\,A_3^2,
  \qquad A_i \in \F_2[x,y],
\end{equation}
and differentiating, with squares derivative-free,
\begin{equation}\label{eq:partials}
  \partial_x q = A_1^2 + y A_3^2, \qquad
  \partial_y q = A_2^2 + x A_3^2, \qquad
  Dq = A_3^2 .
\end{equation}

\begin{lemma}[chain rule for $D$]\label{lem:cocycle}
In characteristic $2$, for all $q, \alpha, \beta \in \F_2[x,y]$,
\[
  D\bigl(q(\alpha,\beta)\bigr)
  = (Dq)(\alpha,\beta)\,
      \bigl(\partial_x\alpha\,\partial_y\beta
            + \partial_y\alpha\,\partial_x\beta\bigr)
  + (\partial_x q)(\alpha,\beta)\, D\alpha
  + (\partial_y q)(\alpha,\beta)\, D\beta .
\]
\end{lemma}

\begin{proof}
Expand $\partial_x\partial_y\bigl(q(\alpha,\beta)\bigr)$ by the chain
and product rules; the terms involving $\partial_x^2 q$ or
$\partial_y^2 q$ vanish by fact (2), and the survivors assemble as
displayed.
\end{proof}

\begin{proposition}\label{prop:Dzero}
If $Dq = 0$ then $xy \notin \Clo(q)$.
\end{proposition}

\begin{proof}
$\{g : Dg = 0\}$ contains the variables and constants and, by
Lemma~\ref{lem:cocycle} with the first term dead, is closed under
substitution into $q$; hence it contains $\Clo(q)$.  But
$D(xy) = 1$.
\end{proof}

Suppose henceforth $Dq \neq 0$, i.e.\ $A_3 \neq 0$; we show
$x + y \notin \Clo(q)$.  Write $\F_2[x{+}y]$ for the subring of
polynomials in $x+y$.  The proof runs through the rigidity property
\begin{quote}
  $(\dagger)$\quad for all $\alpha, \beta \in \F_2[x,y]$: whenever
  $q(\alpha,\beta)$ is a nonconstant element of $\F_2[x{+}y]$, both
  $\alpha, \beta \in \F_2[x{+}y]$.
\end{quote}

\begin{proposition}[master reduction]\label{prop:master}
If $q$ satisfies $(\dagger)$ then $x + y \notin \Clo(q)$.
\end{proposition}

\begin{proof}
We claim every $g \in \Clo(q)$ satisfies
$P(g):\ g \notin \F_2[x{+}y]$ or $g$ constant.  The substitution
$x \leftrightarrow y$ fixes $\F_2[x{+}y]$ pointwise but moves $x$ and
$y$, so the variables satisfy $P$; constants do.  Let
$g = q(\alpha,\beta)$ with $P(\alpha), P(\beta)$ and suppose
$g \in \F_2[x{+}y]$ is nonconstant.  Then $(\dagger)$ forces
$\alpha, \beta \in \F_2[x{+}y]$; $P$ forces them constant; then $g$
is constant --- contradiction.  So $P$ holds on $\Clo(q)$, and
$x + y$ fails $P$.
\end{proof}

Fix an algebraically closed field $K \supseteq \F_2$ of
characteristic $2$ containing an element $\sigma$ transcendental over
$\F_2$; concretely $K = \overline{\F_2(\sigma)}$.  Call $c \in K$
\emph{transcendental} if it satisfies no nonzero polynomial relation
over $\F_2$.

\begin{definition}\label{def:witness}
A \emph{witness} is a pair $\alpha, \beta \in K[t]$, not both
constant, with $q(\alpha, \beta) = c$ constant and $c$
transcendental.
\end{definition}

\begin{proposition}[bridge]\label{prop:bridge}
If no witness exists, then $q$ satisfies $(\dagger)$.
\end{proposition}

\begin{proof}
Let $q(\alpha,\beta) = f(x+y)$ with $f \in \F_2[u]$ nonconstant.
Substitute $(x,y) = (t,\, t + \sigma)$: the right side becomes the
constant $f(\sigma)$, which is transcendental (a nonzero
$\F_2$-relation on $f(\sigma)$ composes with $f$ to one on $\sigma$).
Since no witness exists, $\alpha(t, t{+}\sigma)$ and
$\beta(t, t{+}\sigma)$ are constant in $t$.  Finally, a
$g \in \F_2[x,y]$ constant along the generic line lies in
$\F_2[x{+}y]$: expanding $g(x, x + u) = \sum_i g_i(u)\, x^i$,
constancy gives $g_i(\sigma) = 0$ for $i \geq 1$, hence $g_i = 0$, so
$g = g_0(x+y)$.
\end{proof}

It remains to prove:

\begin{theorem}[no witnesses]\label{thm:descent}
If $Dq \neq 0$, no witness exists.
\end{theorem}

\section{The Frobenius descent}\label{sec:descent}

Throughout, $'$ denotes $d/dt$ and degrees are degrees in $t$.  Over
$K[t]$: squares have zero derivative, and conversely a polynomial
with zero derivative is a square (it is a polynomial in $t^2$, and
the coefficients have square roots since $K$ is perfect).

Set $w := \gcd(A_1, A_2, A_3)$ and $A_i = w B_i$, so the $B_i$ have
no common nonunit factor and
\begin{equation}\label{eq:peel}
  q = A_0^2 + w^2\, h,
  \qquad
  h := x B_1^2 + y B_2^2 + xy B_3^2 ,
\end{equation}
with $B_3 \neq 0$ and partials
$h_x = B_1^2 + y B_3^2$, $h_y = B_2^2 + x B_3^2$, $Dh = B_3^2$.
The peel \eqref{eq:peel} removes the possible curve of critical
points of $q$ along $\{w = 0\}$; Lemma~\ref{lem:keystone} shows none
survives in $h$.  Three lemmas, then the descent.

\begin{lemma}[kill lemma]\label{lem:kill}
Let $(\alpha,\beta)$ be a witness with level $c$, and
$f \in \F_2[x,y]$, $f \neq 0$.  Then $f(\alpha,\beta) \neq 0$ in
$K[t]$.
\end{lemma}

\begin{proof}[Proof sketch]
Suppose $f(\alpha,\beta) = 0$; after swapping variables assume
$\alpha$ nonconstant.  View $f$ and $g := q + z$ as polynomials in
$y$ over $\F_2[x,z]$: since $g$ is monic and linear in $z$ it is
irreducible and cannot divide the $z$-free $f$, so they are coprime,
and a B\'ezout identity with denominators cleared gives
$f\,U + (q + z)\,V = E(x, z) \neq 0$ in $\F_2[x,z][y]$.
Substituting $y \mapsto \beta$, $x \mapsto \alpha$, $z \mapsto c$
kills both summands ($q(\alpha,\beta) + c = 0$), so
$E(\alpha, c) = 0$; expanding $E = \sum_i \rho_i(z)\,x^i$ and using
that $1, \alpha, \alpha^2, \dots$ have strictly increasing degrees in
$t$, every $\rho_i(c) = 0$, and some $\rho_i \neq 0$ exhibits $c$ as
algebraic over $\F_2$ --- contradiction.
\end{proof}

\begin{lemma}[algebraic points]\label{lem:algpoint}
If $u, v \in \F_2[x,y]$ have no common nonunit factor and
$(a,b) \in K^2$ satisfies $u(a,b) = v(a,b) = 0$, then $a, b$ are
algebraic over $\F_2$.
\end{lemma}

\begin{proof}[Proof sketch]
In $\F_2(x)[y]$ a common divisor would descend to $\F_2[x,y]$ by
Gauss's lemma \cite[Ch.~IV]{Lang}, so $u, v$ are coprime there; a
cleared B\'ezout identity gives $su + tv = r(x) \neq 0$, and
evaluation at $(a,b)$ gives $r(a) = 0$.  Symmetrically for $b$.
\end{proof}

\begin{lemma}[keystone]\label{lem:keystone}
Let $B_1, B_2, B_3 \in \F_2[x,y]$ have no common nonunit factor, with
$B_3 \neq 0$, and set $h_x := B_1^2 + y B_3^2$,
$h_y := B_2^2 + x B_3^2$.  Then $h_x$ and $h_y$ have no common
nonunit factor.
\end{lemma}

\begin{proof}
Both are nonzero: applying $\partial_y$ to a relation $h_x = 0$ would
give $B_3^2 = 0$.  Suppose a prime $p$ divides both.

If $p \mid B_3$: subtracting the $B_3^2$-terms, $p$ divides $B_1^2$
and $B_2^2$, hence (primality) all three $B_i$ --- excluded.

If $p \nmid B_3$: write $h_x = p u$ and $h_y = p v$.  Differentiate,
remembering squares are derivative-free:
\[
  0 = \partial_x h_x = p_x u + p\, u_x
  \quad\Longrightarrow\quad p \mid p_x u ;
\]
\[
  B_3^2 = \partial_y h_x = p_y u + p\, u_y
  \quad\Longrightarrow\quad p \nmid u
\]
(else $p \mid B_3^2$, so $p \mid B_3$);
\[
  B_3^2 = \partial_x h_y = p_x v + p\, v_x
  \quad\Longrightarrow\quad p \nmid p_x
\]
(same reason).  So the prime $p$ divides $p_x u$ but neither factor
--- contradiction.
\end{proof}

\begin{proof}[Proof of Theorem~\ref{thm:descent}]
Among all witnesses take one minimizing
$n := \deg\alpha + \deg\beta$ ($\geq 1$, as not both are constant);
fix $r := \sqrt{c}$, transcendental (if $r$ were algebraic, so were
$c = r^2$).  By Lemma~\ref{lem:kill} applied to the nonzero
polynomials $w$, $h_x$, $h_y$, $B_3$, none of
\[
  W := w(\alpha,\beta), \quad
  A := h_x(\alpha,\beta), \quad
  B := h_y(\alpha,\beta), \quad
  \Delta := B_3(\alpha,\beta)^2
\]
vanishes in $K[t]$.

\emph{Peel.}  From \eqref{eq:peel} and $q(\alpha,\beta) = c = r^2$,
using $(u+v)^2 = u^2 + v^2$,
\begin{equation}\label{eq:S}
  W^2 H = S^2,
  \qquad
  H := h(\alpha,\beta), \quad S := r + A_0(\alpha,\beta);
\end{equation}
differentiating ($W^2$ and $S^2$ are squares) gives $W^2 H' = 0$,
so $H' = 0$.

\emph{Engine.}  The chain rule for $h$, with
$\partial_x h_x = \partial_y h_y = 0$ and
$\partial_y h_x = \partial_x h_y = Dh = B_3^2$, gives
\[
  A\alpha' + B\beta' = H' = 0,
  \qquad A' = \Delta\beta', \qquad B' = \Delta\alpha' ,
\]
whence $(AB)' = \Delta(B\beta' + A\alpha') = 0$ and $AB$ is a square
in $K[t]$.

\emph{Fork, branch 1: $A$ and $B$ share a root $\theta \in K$.}  Then
$P := (\alpha(\theta), \beta(\theta))$ is a common zero of
$h_x, h_y$, which by Lemmas~\ref{lem:keystone}
and~\ref{lem:algpoint} has algebraic coordinates.  Evaluating
\eqref{eq:S} at $\theta$, $S(\theta)^2 = w(P)^2\, h(P)$ is algebraic,
hence so is $S(\theta)$, hence so is $r = S(\theta) + A_0(P)$ ---
contradicting transcendence of $r$.

\emph{Fork, branch 2: $A, B$ coprime.}  Coprime factors of a square
are squares (units of $K[t]$ are constants of the algebraically
closed $K$, hence squares), so $A' = B' = 0$; since $\Delta \neq 0$,
the engine identities force $\alpha' = \beta' = 0$, so
$\alpha = \alpha_1^2$, $\beta = \beta_1^2$ with
$\deg\alpha_1 + \deg\beta_1 = n/2$.  Because the coefficients of $q$
lie in $\F_2$, squaring commutes with substitution into $q$:
\[
  q(\alpha_1, \beta_1)^2 = q(\alpha_1^2, \beta_1^2) = c,
\]
so $q(\alpha_1,\beta_1) = r$ and $(\alpha_1, \beta_1)$ is a witness
of total degree $n/2 < n$ --- contradicting minimality.  (If $n = 1$
halving is impossible and this branch cannot occur, so the descent
is well-founded.)
\end{proof}

Theorem~\ref{thm:descent} gives $(\dagger)$ by
Proposition~\ref{prop:bridge}, hence $x+y \notin \Clo(q)$ by
Proposition~\ref{prop:master} when $Dq \neq 0$; with
Proposition~\ref{prop:Dzero} this proves Theorem~\ref{thm:F2}, and
with Lemma~\ref{lem:transfer}, Corollary~\ref{cor:Z}.

\section{From $\F_2$ to the dichotomy}\label{sec:tothedichotomy}

\begin{proposition}\label{prop:perfect}
For every perfect field $k$ of characteristic $2$ and every
$q \in k[x,y]$, at least one of $x+y$, $xy$ lies outside $\Clo(q)$.
\end{proposition}

\begin{proof}
The descent runs over $k$ with one change: coefficients of $q$ are no
longer fixed by squaring, so in the halving branch
$q(\alpha_1^2, \beta_1^2) = q^{(1/2)}(\alpha_1, \beta_1)^2$, where
$q^{(1/2)}$ applies the inverse Frobenius to each coefficient
(perfectness enters here); the halved pair is a witness for
$q^{(1/2)}$, and since the induction is on witness degree with $q$
universally quantified inside, the twist is absorbed.  Parity
decomposition, bridge, and master reduction are unaffected.
\end{proof}

\begin{proof}[Proof of Theorem~\ref{thm:char2}]
$k$ embeds into its perfection $k^{1/2^\infty}$; by
Lemma~\ref{lem:transfer}, membership of both $x+y$ and $xy$ in
$\Clo(q)$ over $k$ would persist over the perfection, contradicting
Proposition~\ref{prop:perfect}.
\end{proof}

\begin{proof}[Proof of Theorem~\ref{thm:dichotomy}]
If $2$ is a unit, $x^2 - y$ is complete
(Theorem~\ref{thm:converse}).  If not, $2$ lies in a maximal ideal
$\mathfrak m$; a complete operation would push forward along
$R \to R/\mathfrak m$ (Lemma~\ref{lem:transfer}) to one over the
residue field, of characteristic $2$ --- contradicting
Theorem~\ref{thm:char2}.
\end{proof}

\medskip

All results have been machine-verified in Lean~4 against Mathlib
(axioms: propositional extensionality, quotient soundness, choice);
the development accompanies the paper.

Three questions remain open.  Which degrees $\geq 3$ admit complete
operations over $\Q$ ($x^3 - y$ is a candidate; over $\F_3$ its clone
consists of additive-plus-constant polynomials, so it is not complete
there)?  Which $p$ are complete when $2$ is a unit ($xy + 1$ and
$x^2 + y$ are not, by a characteristic-zero instance of $(\dagger)$
and by a positivity obstruction, respectively)?  Is membership in
$\Clo(p)$ --- or completeness --- decidable?  The identity
$q(f, q(f,g)) = g$ for $q = x^2 + y$ over $\F_2$ shows substitution
admits unbounded cancellation, so bounded search proves nothing.

\begin{thebibliography}{9}

\bibitem{Lang} S.~Lang, \emph{Algebra}, rev.\ 3rd ed., Graduate Texts
in Mathematics 211, Springer, New York, 2002.

\bibitem{Sheffer} H.~M.~Sheffer, \emph{A set of five independent
postulates for Boolean algebras, with application to logical
constants}, Trans.\ Amer.\ Math.\ Soc.\ 14 (1913), 481--488.

\end{thebibliography}

\end{document}
